%(BEGIN_QUESTION)
% Copyright 2006, Tony R. Kuphaldt, released under the Creative Commons Attribution License (v 1.0)
% This means you may do almost anything with this work of mine, so long as you give me proper credit

En instrumenttekniker vil lage en temperaturreferanse for en termoelementtransmitter ved å fryse vann, siden frysepunktet for vann ved havnivå er 32$^{o}$ F (0$^{o}$ C). Han setter termoelementet i en kopp vann, plasserer koppen og termoelementet i en fryser til vannet er helt frosset, tar deretter koppen ut av fryseren og kobler termoelementet til transmitteren for kalibrering.

Hva er feil i denne prosedyren? Hva må gjøres annerledes for å sikre en referansetemperatur på 32$^{o}$ F (0$^{o}$ C) ved termoelementspissen?

\vskip 20pt \vbox{\hrule \hbox{\strut \vrule{} {\bf Forslag til sokratisk diskusjon} \vrule} \hrule}

\begin{itemize}
\item{} Bruk et \textit{fasediagram} for å identifisere punkter der H$_{2}$O holder svært stabil temperatur selv ved endringer i omgivende trykk.
\item{} Forklar hvordan vannets \textit{trippelpunkt} brukes i metrologilaboratorier for å opprettholde både stabil temperatur og stabilt trykk i kalibreringssammenheng.
\end{itemize}

\underbar{file i00355}
%(END_QUESTION)





%(BEGIN_ANSWER)

%(END_ANSWER)





%(BEGIN_NOTES)

Det teknikeren trenger, er en \textit{blanding av is og vann} for å garantere stabilitet ved frysetemperaturen.

\vskip 10pt

I stedet for å sette en kopp vann i fryseren anbefales det å legge isbiter i en kopp vann og røre blandingen for jevn temperatur. Fortsett å fylle på med is etter hvert som den smelter, så holder blandingen konstant 32$^{o}$ F (0$^{o}$ C) uavhengig av romtemperaturen.

%INDEX% Calibration, standard: using an ice bath properly

%(END_NOTES)
